\documentclass[../notes.tex]{subfiles}
\begin{document}

\chapter{Conventions and notation}
\markboth{Conventions and notation}{}

These conventions hold in every chapter. When a chapter needs new notation, it is added here first, so the whole book stays consistent. Most confusion in robotics comes from mismatched conventions, not from hard mathematics.

\section*{Frames and coordinates}
\begin{itemize}
\item Frames are written $\{A\},\{B\},\dots$; reserved names: $\{W\}$ world, $\{B\}$ robot base/body, $\{C\}$ camera, $\{I\}$ IMU, $\{E\}$ end-effector. $\pre{A}{p}$ = coordinates of point/vector $p$ expressed in frame $\{A\}$.
\item $R_{AB}\in\SO$ is the orientation of $\{B\}$ \emph{relative to} $\{A\}$. As a matrix it converts coordinates \emph{from} $B$ \emph{to} $A$: $\pre{A}{p}=R_{AB}\,\pre{B}{p}$. Inner indices cancel: $R_{AC}=R_{AB}R_{BC}$.
\item Poses: $T_{AB}\in\SE$, with $t_{AB}$ the origin of $\{B\}$ written in $\{A\}$; $\pre{A}{p}=R_{AB}\pre{B}{p}+t_{AB}$.
\item ROS body convention: $x$ forward, $y$ left, $z$ up. Camera optical convention: $z$ forward, $x$ right, $y$ down.
\end{itemize}

\section*{Rotations and their parameterisations}
\begin{itemize}
\item $\skw{\omega}$: the $3\times3$ skew matrix with $\skw{\omega}v=\omega\times v$. $\Exp:\RR^3\to\SO$ and $\Log:\SO\to\RR^3$ map between rotation vectors and matrices; $e^{(\cdot)}$ is the ordinary matrix exponential.
\item Euler angles: aerospace/ROS ``RPY'' $R=R_z(\psi)R_y(\theta)R_x(\phi)$: yaw $\psi$, pitch $\theta$, roll $\phi$.
\item Quaternions: Hamilton ($ij=k$), scalar first, $q=(w,x,y,z)=(w,u)$, product $\otimes$. ROS messages, Eigen's storage and SciPy output use $(x,y,z,w)$.
\item Angular velocity: $\dot R=\skw{\omega_s}R=R\skw{\omega_b}$ ($s$ = spatial/world, $b$ = body).
\end{itemize}

\section*{Planar motion (chapter~\ref{ch:pl})}
\begin{itemize}
\item $J=\bmat{0&-1\\1&0}$ is the planar rotation generator ($J^2=-I$; $Jp$ is $p$ turned by $90^\circ$). $R(\theta)=e^{\theta J}\in\SOtwo$. (Chapter~\ref{ch:rot} uses $J$ for Jacobians; the context always makes clear which is meant.)
\item Planar poses $T=\bmat{R(\theta)&t\\0&1}\in\SEtwo$. Planar twists are linear-first like 3-D ones: $S=(v_x,v_y,\omega)$, $[S]=\bmat{\omega J&v\\0&0}\in\setwo$; $\vexa$ is the inverse of $[\cdot]$. A unit twist $\hat S$ is scaled by a motion amount $\lambda$ (an angle or a distance): $T=\exp(\lambda[\hat S])$.
\end{itemize}

\section*{Lie groups (chapter~\ref{ch:lie})}
\begin{itemize}
\item A generic group element is $\mathcal X$, the identity $\mathcal E$; $m$ = degrees of freedom. Tangent vectors $\tau\in\RR^m$, hat $\tau^\wedge$ into the Lie algebra, vee back. $\Exp(\tau)=\exp(\tau^\wedge)$, $\Log$ its inverse near $\mathcal E$.
\item Right (local) plus/minus by default: $\mathcal X\oplus\tau=\mathcal X\Exp(\tau)$, $\mathcal Y\ominus\mathcal X=\Log(\mathcal X^{-1}\mathcal Y)$. Left (global) versions are stated explicitly.
\item $\Ad_{\mathcal X}$: adjoint matrix. $J_r(\tau)$, $J_l(\tau)$: right and left Jacobians of $\Exp$. Jacobians of functions are right Jacobians unless stated. Notation follows Solà, Deray and Atchuthan (2018).
\end{itemize}

\section*{ROS 2 and tf2 (section~\ref{sec:rot:tf2})}
\begin{itemize}
\item A \texttt{TransformStamped} with \texttt{header.frame\_id=A}, \texttt{child\_frame\_id=B} is $T_{AB}$: the pose of $B$ in $A$, mapping $B$-coordinates to $A$-coordinates. \texttt{lookup\_transform(target, source)} returns $T_{\text{target},\text{source}}$; \texttt{tf2\_echo A B} prints $T_{AB}$.
\item ROS quaternions are $(x,y,z,w)$; this book writes $(w,x,y,z)$. ROS twists (\texttt{geometry\_msgs/Twist}) are linear-first, like ours.
\item REP~103: body frames $x$ forward, $y$ left, $z$ up; \texttt{\_optical} frames $z$ forward, $x$ right, $y$ down. REP~105: \texttt{map}$\to$\texttt{odom}$\to$\texttt{base\_link}.
\end{itemize}

\section*{Rigid-body motion}
\begin{itemize}
\item Twists are ordered linear-first, $\xi=(v,\omega)\in\RR^6$, as in Murray--Li--Sastry, Barfoot and Sophus. Lynch \& Park (\emph{Modern Robotics}) and GTSAM use $(\omega,v)$; the adjoint blocks swap accordingly.
\item $[\xi]=\bmat{\skw{\omega}&v\\0&0}\in\se$; $\Ad_T=\bmat{R&\skw{t}R\\0&R}$.
\item Perturbations are on the right (body frame) unless stated: $R\Exp(\delta)$, $T\Exp(\delta)$.
\end{itemize}

\section*{General}
\begin{itemize}
\item Vectors are columns. $c_\theta=\cos\theta$, $s_\theta=\sin\theta$. Angles in radians in formulas, degrees in examples.
\item Box types: \textbf{Key idea} (the essence), \textbf{Derivation} (full argument), \textbf{Worked example} (applied, numbers verified in code), \textbf{Pitfall} (what goes wrong in practice), plus plain note boxes.
\end{itemize}

\end{document}
